\backhead{OBJECTIVE TYPE QUESTIONS}{khred}
\begin{exlist}
\item A two-level full factorial of six factors has\opt{12}{36}{64}{720} runs
\item In leave-one-soil-out validation of ten soils each fit uses\opt{1}{5}{9}{10} soils
\item A bootstrap draws soils\opt{without replacement}{with replacement}{in order}{once each}
\item A Latin hypercube samples each factor\opt{at two levels}{once in each stratum}{at random only}{at its mean}
\item An overfitted law has\opt{a low LOSO error}{a low in-sample error and a high LOSO error}{no parameters}{a high in-sample error}
\end{exlist}
\objkey{1.\ (c),\quad 2.\ (c),\quad 3.\ (b),\quad 4.\ (b),\quad 5.\ (b)}
\dykband{EXERCISES IN SWAFL}
{\small Open SWAFL (desktop or online) with the benchmark project. The chapter of the \href{https://github.com/halilunaziru73-creator/Coupled-Water-Fluxes-Pore-Scale-Dynamics-Precision-Irrigation-Mulching-Feedback-Loops/raw/main/Software/v1.1.0/SWAFL_Desktop/docs/manual_v1.1.0/SWAFL_Manual_v1.1.0.pdf}{SWAFL manual} that explains each task is given in brackets; record your values in a table and compare them with what is expected (\emph{Check}).}\par\vspace{3pt}
\begin{exlist}
\item {\color{khblue}[Manual, Chapter 28]} Design a Daisy experiment with a Latin hypercube of 20 runs and inspect the plan.\par\noindent{\small\itshape Check: each factor is sampled in 20 strata.}
\item {\color{khblue}[Manual, Chapter 29]} Open the best combinations and choose one for your field.\par\noindent{\small\itshape Check: the plan of the combination is shown.}
\item {\color{khblue}[Manual, Chapter 30]} Read the interactions of an experiment and identify the strongest.\par\noindent{\small\itshape Check: usually between water and nitrogen.}
\item {\color{khblue}[Manual, Chapter 18]} Use the search for the most efficient inputs and compare with the best combination of Daisy.\par\noindent{\small\itshape Check: the two agree in direction.}
\end{exlist}
